\section{Collation Findings and Contested Readings}
\label{sec:collation}

\subsection{Method and its limits}

Two witnesses control this reference: the promulgated Latin of the code in \latin{Acta
Apostolicae Sedis} 52 (1960), read in the Holy See's own digitisation, and the code as reprinted
at pages XIII--XXXVI of the 1962 Vatican typical edition of the Missal, read in the facsimile
published by the Church Music Association of America. Both digitisations carry optical text
layers whose readings are discovery aids only. Every reading reported as a divergence below was
read in a rendered page image at 260 to 450 dots per inch, on both sides, and is not an artefact
of optical recognition. No third exemplar of either printing was consulted, so a defect of the
particular copy scanned cannot be excluded; the findings are stated as readings of the exact
bytes identified in the source records, not as claims about every printed copy.

\subsection{Four divergences between the Acta and the Missal}

\begin{center}
\small
\renewcommand{\arraystretch}{1.14}
\begin{tabular}{@{}>{\raggedright\arraybackslash}p{0.16\linewidth} >{\raggedright\arraybackslash}p{0.30\linewidth} >{\raggedright\arraybackslash}p{0.22\linewidth} >{\raggedright\arraybackslash}p{0.24\linewidth}@{}}
\toprule
\textbf{Locus} & \textbf{Acta 52 (1960)} & \textbf{Missal 1962} & \textbf{Assessment} \\
\midrule
\rn{14}, p.~598 & Numbered \textbf{11} & Numbered \textbf{14} & Acta misprint; nn.~11 and 14 cannot both be 11, and the surrounding numbers 12, 13, 15 are correct \\
\midrule
\rnn{15} and 16, p.~598 & \latin{servetur norma, quae nn.\ \textbf{101-105} traditur}, twice & \latin{nn.\ \textbf{104-105}} & Missal reading is substantively right: the concurrence norms are \rnn{103--105}; \rnn{101--102} govern reposition \\
\midrule
\rn{18}\,d, p.~599 & \latin{et \textbf{XVII}, quae inscribitur VI} & \latin{et \textbf{XXVII}, quae inscribitur VI} & Acta misprint; a seventeenth place cannot follow a twenty-sixth, and case \textit{d} concerns years of 28 Sundays \\
\midrule
Missal decree, 23 June 1962 & Motu proprio dated \latin{die 25 mensis Iulii, anno 1960}; decree dated 26 July & Decree recites \latin{Motu proprio ``Rubricarum instructum'' diei \textbf{23 iulii} anno 1960 \dots\ posteroque die \dots\ promulgato} & Missal misprint; the same volume reprints the motu proprio later in its own front matter with its own date, \latin{die 25 mensis iulii, anno 1960}, and 23 July plus one day is not 26 July, so the decree is inconsistent with the book that carries it \\
\bottomrule
\end{tabular}
\end{center}

The volume's own \latin{Corrigenda} notice, printed at the end of \latin{Acta} 52, records a
single correction, to page 179. It does not reach pages 598 or 599. The two Acta misprints
therefore stood uncorrected in the volume as issued.

None of the four affects the substance of the law as this reference states it, because in each
case one reading is impossible and the other is not. But three of them fall in the calendar's
governing text, and a reference that cited only the promulgating gazette would print an
arithmetically impossible ordinal in the rule that governs the resumption of the Epiphany
Sundays. That is a practical reason for a calendar reference to collate the typical edition
against the Acta rather than treating either as sufficient alone.

\subsection{A misprint inside the Missal's own table}

The \latin{Tabella temporaria festorum mobilium} printed among the paschal tables of the 1962
typical edition gives, for each year from 1960 onward, the number of Sundays after Pentecost. For
1996 the cell reads \textbf{16}.

\begin{itemize}[leftmargin=1.3em, itemsep=0.25em]
\item Sixteen is outside the possible range. The same book's perpetual \latin{Tabula paschalis
antiqua reformata}, printed in the same run of tables, shows the number ranging from 28 at the
earliest Easters to 23 at the latest, and no combination of epact and dominical letter yields
fewer than 23.
\item The perpetual table itself gives the right value for the 1996 combination. Its row for
Easter on 7 April with Pentecost on 26 May, Corpus Christi on 6 June, and the First Sunday of
Advent on 1 December — which is exactly the 1996 line of the temporary table — prints
\textbf{26} in the Sundays-after-Pentecost column.
\item Computation agrees with the perpetual table: 26 May to 1 December is 189 days, that is
27 weeks, so $P = 27 - 1 = 26$.
\end{itemize}

The 1996 cell is therefore a printing or scanning defect in the temporary table, contradicted by
the perpetual table in the same book. All fifty-one other rows of that page
were recomputed and agree with the printed values in every column.

This is the kind of finding a reference should report rather than silently correct, because a
reader using the Missal's table directly — as the book intends — would otherwise reach an
impossible result and have no way to know why.

\subsection{A changed rule that looks like a constant}

The Missal's front-matter note \latin{De anno et eius partibus} fixes the Ember Days as the
Wednesday, Friday, and Saturday after the third Sunday of Advent, after the first Sunday of Lent,
after Pentecost Sunday, and \latin{post dominicam tertiam septembris}, after the third Sunday of
September. The corresponding note in the pre-1955 comparison witness reads \latin{post Festum
Exaltationis sanctae Crucis}, after the feast of the Exaltation of the Holy Cross.

The two rules are not equivalent. Under the 1962 wording and \rn{19}, the third Sunday of
September is the Sunday falling from 15 to 21 September, and the Ember Wednesday falls from 18 to
24 September. Under the older wording, the Ember Wednesday is the first Wednesday after
14 September and can fall as early as 15 September. The difference is up to a full week.

The new wording is not a peculiarity of the Vatican printing. A third witness registered in the
repository's source library — the Benziger Brothers \latin{editio iuxta typicam} of New York,
1962, whose approval leaf recites that the Sacred Congregation of Rites declared it
\latin{iuxta typicam} on 21 October 1961 and which therefore carries the 1960 code — prints the
same note in the same place with \latin{post dominicam III septembris}. Two independently
produced books conformed to the 1960 code agree against the earlier Benziger printing, so the
divergence is between states of the rule and not between compositors.

\subsection{The comparison witness, identified}
\label{sec:witness}

The Internet Archive item that carries this witness supplies only an uploader-supplied title
naming 1947 and no publisher, date, or rights metadata, and an earlier draft of this reference
accordingly treated the witness as of wholly unresolved identity. Reading the scan's own opening
leaves corrects that. The book identifies itself on its title page as \latin{Missale Romanum ex
decreto Sacrosancti Concilii Tridentini restitutum, S. Pii V Pontificis Maximi iussu editum,
aliorum Pontificum cura recognitum, a Pio X reformatum, Benedicti XV auctoritate vulgatum},
\latin{editio IV iuxta typicam Vaticanam}, printed at New York by Benziger Brothers, Inc.,
\latin{Summi Pontificis et Sacrae Rituum Congregationis typographorum}. It carries a copyright
line dated 1942 and 1944; an approval leaf subscribed by Francis Joseph, Archbishop of New York,
given at New York on 8 December 1942, certifying that the book was found to agree with the
typical edition and with the additions and variations of the most recent decree of the Sacred
Congregation of Rites of 9 January 1942; and a reprint of the Congregation's decree of
25 July 1920 declaring the Vatican edition prepared under \latin{Divino afflatu} typical.

The witness is therefore a licensed publisher's edition conformed to the Vatican typical edition
of 1920, in a setting approved for New York at the end of 1942 — not the typical edition itself,
and not of established printing year, since nothing in the scan corroborates the uploader's
``1947''. That is a materially stronger position than an unidentified digitisation, and it
supports more than the two observations the earlier draft allowed. Four readings are taken from
it in this reference, each cited at its place: the older September Ember wording above; the
resumption rubric printed after the Twenty-third Sunday after Pentecost, whose wording matches
the 1962 Missal's apart from its internal page cross-reference and the grade
\latin{Semiduplex} given to the resumed formularies; the pre-1960 rule at
\latin{Additiones et Variationes}, \latin{De occurrentia et de translatione festorum}, \rn{6},
by which an impeded Sunday's Mass \latin{resumitur prima infra hebdomadam die} on a free weekday
of the following week; and the pre-1960 ceiling on orations, which required that they
\latin{septenarium numerum non excedant, atque imparem praeterea numerum} servent — not exceed
the number seven, and keep besides an odd number.

What the witness cannot supply is the pre-1955 precedence machinery. Under the older
arrangement, as under the new one at \rn{5} of \latin{Rubricarum instructum}, the tables of
occurrence and concurrence belonged to the Breviary; the Missal's own general rubrics and
\latin{Additiones et Variationes} treat occurrence in the comparative language of an
\latin{Officium nobilius}, a nobler Office, and contain no table and no exclusivity clause. That
is a bounded negative result on the Missal side only, and the collected list below keeps it
distinct from the Breviary-side question that remains open.

Finally, this reference did not locate the act that changed the September Ember rule, which may
be the general decree of 23 March 1955, the 1960 code's \latin{Variationes}, or an editorial
decision of the 1962 typical edition. The \latin{Variationes in Calendario} read for this
document do not mention the Ember Days. The change is reported as observed; its date and
instrument are not established here.

\subsection{Contested and unresolved points, collected}

\begin{enumerate}[leftmargin=1.6em, itemsep=0.35em]
\item \textbf{The year of twenty-three Sundays after Pentecost.} Neither the code nor the Missal
legislates the case. Section~\ref{sec:resumed} gives both readings, their textual support, and
the witness that would settle the question. Unresolved.
\item \textbf{Which II class feasts are \latin{festa Domini}.} The distinction is operative at
lines 14 and 16 of the table of precedence and at \rn{16}\,a, and it decides whether a II class
feast absorbs a II class Sunday or is commemorated by it. The 1962 calendar pages do not mark
the category, and this reference did not establish a controlling list. The concrete instance of
9 November 2008 is left open in Section~\ref{sec:years}. Unresolved.
\item \textbf{The date instrument for the September Ember Days.} Observed as changed; the act
that changed it is not identified here. Unresolved.
\item \textbf{Whether the pre-1960 rubrics contained an exclusivity clause.} Bearing directly on
the judgement in the Project Synthesis. Partly tested: the identified Benziger \latin{editio IV}
described in Section~\ref{sec:witness} carries no such clause on the Missal side and resolves
occurrence by the comparative \latin{Officium nobilius}. The pre-1955 tables of occurrence and
concurrence belonged to the Breviary, which that witness is not, and no pre-1955 Breviary was
read. Unresolved on the Breviary side, and flagged in the Project Synthesis as the decisive test.
\item \textbf{The Breviary tables of occurrence and concurrence.} These are printed in the Acta
after the code and belong to the Office; they are described in Section~\ref{sec:precedence} but
not reproduced, and their detail was not collated. Out of scope rather than contested.
\end{enumerate}
